% Generated draft: requires mathematical and source review.
\documentclass[11pt]{article}
\usepackage[margin=25mm]{geometry}
\usepackage{fontspec}
\setmainfont{Latin Modern Roman}
\usepackage{amsmath,amssymb,mathtools,unicode-math}
\setmathfont{Latin Modern Math}
\usepackage{enumitem,needspace,xurl,hyperref}
\setlength{\parindent}{0pt}
\setlength{\parskip}{5pt}
\setlength{\emergencystretch}{3em}
\newcommand{\N}{\mathbb{N}}
\newcommand{\Z}{\mathbb{Z}}
\newcommand{\Q}{\mathbb{Q}}
\newcommand{\R}{\mathbb{R}}
\newcommand{\C}{\mathbb{C}}
\newcommand{\F}{\mathbb{F}}
\newcommand{\A}{\mathbb{A}}
\newcommand{\PP}{\mathbb{P}}
\newcommand{\E}{\mathbb{E}}
\newcommand{\eps}{\varepsilon}
\newcommand{\dd}{\,\mathrm{d}}
\begin{document}
\section*{UnsolvedMath ID: 30006255}
\textbf{Source problem code:} catalogue entry
\subsection{Definitions and mathematical statement}
Let \(k\) denote the ground field intended by the source, and let \(X\) be an affine \(k\)-variety. The supplied record does not specify hypotheses on \(k\), or whether varieties must be irreducible or reduced. These conventions require confirmation; no algebraic-closure or characteristic assumption is imposed here.

Write \(\operatorname{Aut}(X)\) for the group of algebraic automorphisms of \(X\), interpreted with its usual ind-group structure. Closedness, closure, connectedness, and the identity component \(\operatorname{Aut}(X)^\circ\) are understood in the corresponding ind-group topology. For a closed subgroup \(G\), write \(\operatorname{Lie}(G)=T_eG\) for its tangent Lie algebra at the identity, viewed inside the Lie algebra \(\operatorname{Vec}(X)=\operatorname{Der}_k(k[X])\) of algebraic vector fields. The intended ground-field, point, and ind-group conventions must be checked against the report.

For this draft, an element of \(\operatorname{Aut}(X)\) is called algebraic if it belongs to an algebraic subgroup of \(\operatorname{Aut}(X)\). A subgroup \(G\) is called nested if it admits an increasing exhaustion
\[
G=\bigcup_{i\geq 1}G_i,\qquad G_i\subseteq G_{i+1},
\]
by algebraic subgroups, with the inclusions compatible with its ind-group structure. The record does not give these definitions explicitly, including whether any connectedness condition is part of its nestedness convention.

The corrected statement in the supplied record asks whether the following assertions hold universally under the intended conventions.
\begin{enumerate}
\item For every connected closed subgroup \(G\subseteq\operatorname{Aut}(X)\) and every proper closed normal subgroup \(N\triangleleft G\),
\[
\operatorname{Lie}(N)\subsetneq\operatorname{Lie}(G).
\]
\item For all closed subgroups \(G,H\subseteq\operatorname{Aut}(X)\), if \(G\) is connected, every element of \(G\) is algebraic, and
\[
\operatorname{Lie}(G)\subseteq\operatorname{Lie}(H),
\]
then \(G\subseteq H\). These hypotheses come from the corrected statement: the supplied original transcription omits them and merely asks whether \(G\subseteq H\).
\item Every closed subgroup of \(\operatorname{Aut}(X)\) consisting entirely of algebraic elements is nested. More generally, for every subgroup \(G\subseteq\operatorname{Aut}(X)\) consisting entirely of algebraic elements, its closure \(\overline G\) is nested. No additional assertion that every element of \(\overline G\) is algebraic is assumed. The generalization to arbitrary, possibly nonclosed \(G\) follows the corrected statement; the original transcription first assumes \(G\) closed and then asks about its closure, leaving the intended distinction unclear.
\item If \(\operatorname{Aut}(X)^\circ\neq\{e\}\), then there exists a positive-dimensional algebraic subgroup
\[
A\subseteq\operatorname{Aut}(X)^\circ.
\]
The original wording describes the premise as nondiscreteness and asks for a nonfinite algebraic subgroup, further describing this as a copy of \(k^*\) or \(k^+\), conventionally represented by \(\mathbb G_m\) or \(\mathbb G_a\). Those equivalences require ground-field and structural hypotheses absent from the record, and are retained as source qualifications rather than asserted here without qualification.
\item In the notation of the source, put
\[
S_n=\operatorname{SAut}(\mathbb A_k^n),\qquad
T_n=\langle\operatorname{SAff}_n,\operatorname{SJonq}_n\rangle\subseteq S_n.
\]
For a subgroup \(T\subseteq S\), its normalizer is
\[
N_S(T)=\{s\in S:sTs^{-1}=T\}.
\]
The question is whether
\[
\overline{N_{S_n}(T_n)}^{\,S_n}=S_n,
\]
where the closure is taken in the ambient ind-group \(S_n\). The definitions of \(\operatorname{SAut}\), \(\operatorname{SAff}_n\), and \(\operatorname{SJonq}_n\), and the intended range of \(n\), are not supplied. Thus this question preserves the exact subgroup notation but remains conditional on recovering those conventions.
\end{enumerate}

The surrounding source passage supplies relevant qualifications. It attributes to Kraft--Zaidenberg (2014), with reference [KZ24a], a result that a subgroup generated by algebraic subgroups \(\{G_j\}_{j\in J}\) is an algebraic group if and only if the Lie algebras \(\operatorname{Lie}(G_j)\) generate a finite-dimensional Lie algebra of vector fields. It also attributes to [FK18] an example of a connected closed subgroup \(G\subseteq\operatorname{Aut}(\mathbb A^3)\) with a proper closed subgroup \(H\subsetneq G\) satisfying \(\operatorname{Lie}(H)=\operatorname{Lie}(G)\), and says this cannot occur when \(G\) is nested. The passage does not say that this \(H\) is normal, so it does not supply a counterexample to question (1). These are supplied contextual claims, with any omitted hypotheses unresolved, rather than independently verified results or additional problems to prove.
\subsection{Short English statement}
The record asks five questions about how closed automorphism subgroups relate to their Lie algebras, whether groups of algebraic elements are nested, whether a nontrivial identity component contains a positive-dimensional algebraic subgroup, and whether the normalizer of the stated tame subgroup is dense. The supplied curation labels all five questions open. Ground-field assumptions, several subgroup conventions, and discrepancies between the original transcription and corrected statement still require mathematical review.
\subsection{Sources}
\begin{itemize}
\item Supplied local UnsolvedMath registry record, permanent ID \(30006255\), title \emph{Open Problems on Ind-Groups of Automorphisms}, source code \texttt{catalogue entry}, in the Oberwolfach Reports Open Problems collection. The record gives proposed year 2025 and names no individual proposer.
\item Supplied report citation: \emph{Algebraic Groups} (2025), Oberwolfach Reports, pp. 941--1010, DOI 10.4171/owr/2025/20. \url{https://doi.org/10.4171/owr/2025/20}. The background identifies Oberwolfach Reports as reports of mathematical workshops and their open-problem discussions.
\item Supplied literature source, labelled ``Kraft survey, open problems'': \url{https://doi.org/10.4171/OWR/2025/20}.
\item Supplied literature source, labelled ``Kraft--Zaidenberg, nested automorphism groups'': \url{https://arxiv.org/abs/2312.08359}.
\item The source passage additionally cites [PR23] and [PR24] for the nestedness questions, [KZ24b] for an overview of relations between subgroups and Lie algebras, [KZ24a] for the result attributed to Kraft--Zaidenberg (2014), and [FK18] for the equal-Lie-algebra example. Full bibliographic entries and the identification of these labels with the supplied links are not provided.
\item Literature assessment dated exactly \texttt{2026-08-22}: the supplied curation labels the five questions open and states that recent structural results do not settle them in general. This is supplied dated curation, explicitly described as literature triage requiring independent verification; it has not been independently verified for this draft.
\item The registry labels the statement \texttt{corrected\_verified} and records that the official report was checked and its notation repaired. This is a supplied source-verification claim, not a claim that the report was visited or checked while preparing this draft. The supplied record was updated on \texttt{2026-09-16T12:06:53Z}.
\end{itemize}
\end{document}
